\section{The physical boundary symmetry of a CZX rectangle}%
\label{sec:peps_czx_rectangle_boundary}

Let $R=[x_0,x_0+w)\times[y_0,y_0+h)$ be a coordinate rectangle in a
$W\times H$ torus. Assume $W,H\geq3$, $w,h>0$, $w<W$, $h<H$, and
$x_0+w\leq W$, $y_0+h\leq H$. Write $P=2w+2h$.
The following construction uses the actual open-region map $A_R$, summing
only internal bonds and leaving crossing bonds free. It identifies the
plaquette spins in the boundary discussion of \cite{ChenLiuWen2011CZX},
source lines 330--345, with the image of this physical region map.

\begin{definition}[Clockwise native rectangle boundary]%
    \label{def:czx_rectangle_perimeter}
    \lean{TNLean.PEPS.rectanglePerimeterCorner,
        TNLean.PEPS.rectanglePerimeterEnd,
        TNLean.PEPS.rectanglePerimeterReversed,
        TNLean.PEPS.torusRectanglePerimeterCode,
        TNLean.PEPS.torusRectanglePerimeterEdge,
        TNLean.PEPS.torusRectanglePerimeterSite,
        TNLean.PEPS.torusRectanglePerimeterLeg,
        TNLean.PEPS.torusRectanglePerimeterCorner}
    \leanok
    \uses{def:peps_torus_rectangular_grid}
    Number the crossing bonds clockwise, first across the top, then down
    the right, across the bottom, and up the left. The corresponding first
    plaquette corners are
    $(i,h)$, $(w,h-j)$, $(w-k,0)$, and $(0,\ell)$ on the four sides.
    Reverse the native pair order on the bottom and left.
\end{definition}

\begin{theorem}[Native perimeter and distinct plaquette corners]%
    \label{thm:czx_rectangle_perimeter}
    \lean{TNLean.PEPS.even_rectanglePerimeter,
        TNLean.PEPS.rectanglePerimeterEnd_eq_corner_add_one,
        TNLean.PEPS.torusRectanglePerimeterEdge_boundary,
        TNLean.PEPS.torusRectanglePerimeterEdge_injective,
        TNLean.PEPS.torusRectangleBoundaryEquiv,
        TNLean.PEPS.torusRectangleBoundaryEquiv_apply_val,
        TNLean.PEPS.torusRectanglePerimeterSite_mem,
        TNLean.PEPS.torusRectanglePerimeterLeg_val,
        TNLean.PEPS.rectanglePerimeterCorner_le,
        TNLean.PEPS.rectanglePerimeterCorner_injective,
        TNLean.PEPS.torusRectanglePerimeterCorner_injective}
    \leanok
    \uses{def:czx_rectangle_perimeter}
    This numbering is a bijection from $\{0,\ldots,P-1\}$ to the actual
    crossing bonds. Its cyclic successor is ordinary addition in
    $\mathbb Z/P\mathbb Z$. Every bond has the specified endpoint in $R$,
    the perimeter corners are distinct in the torus, and $P$ is even.
    The assertions include rectangles meeting a coordinate seam.
\end{theorem}
\begin{proof}\leanok
    \uses{lem:peps_torus_rectangle_boundary_card}
    The four native edge families are disjoint and injective. The existing
    crossing-bond cardinality theorem gives surjectivity. Proper side
    lengths make the closed corner grid injective modulo each torus period.
\end{proof}

\begin{theorem}[Exact coefficients and disjoint boundary columns]%
    \label{thm:czx_open_region_coefficients}
    \lean{TNLean.PEPS.czx_regionIncidentConfig_unique,
        TNLean.PEPS.openRegionWeight_czxPEPS,
        TNLean.PEPS.openRegionWeight_czxPEPS_disjoint,
        TNLean.PEPS.openRegionCoefficientMatrix_czxPEPS_gram_offDiagonal}
    \leanok
    \uses{def:peps_open_region_contraction,def:peps_czx_tensor}
    For any finite region, a physical configuration determines at most one
    nonzero incident-bond assignment. Its open-region coefficient is one
    precisely when this assignment exists and has the specified boundary;
    otherwise it is zero. Different boundary columns have disjoint physical
    supports, so the region Gram matrix is diagonal in the native boundary
    basis. No labels on edges wholly outside the region are summed.
\end{theorem}
\begin{proof}\leanok
    The printed CZX tensor fixes all four incident labels from the four
    physical qubits. Every incident edge touches a region vertex, so this
    determines the entire assignment. The contraction sum has at most one
    nonzero summand, whose value is one.
\end{proof}

\begin{definition}[Actual crossing-bond symmetry]%
    \label{def:czx_region_boundary_symmetry}
    \lean{TNLean.PEPS.czxRegionBoundaryFlip,
        TNLean.PEPS.czxRegionBoundaryPhase}
    \leanok
    \uses{def:peps_czx_leg_operator,def:peps_open_region_contraction}
    For a native boundary configuration $\mu$, let $f_R\mu$ flip both bits
    on each crossing bond and let
    $\phi_R(\mu)=\prod_{e\in\partial R}(-1)^{a_eb_e}$, where
    $\mu_e=(a_e,b_e)$.
\end{definition}

\begin{theorem}[Physical region pull-through]%
    \label{thm:czx_region_physical_pullthrough}
    \lean{TNLean.PEPS.czxLegPhase_swap,
        TNLean.PEPS.czxLegFlip_swap,
        TNLean.PEPS.regionPhysicalMap_czxOnSite_openRegionWeight,
        TNLean.PEPS.regionPhysicalProductMatrix_czxOnSite_intertwines}
    \leanok
    \uses{def:czx_region_boundary_symmetry,def:peps_czx_on_site}
    Let $U_R$ apply the CZX on-site symmetry at each site of $R$. Then
    \[
        U_RA_R|\mu\rangle=\phi_R(\mu)A_R|f_R\mu\rangle.
    \]
    This identity holds for the actual open contraction of every finite
    region. Reversal of a native bond pair preserves its phase and commutes
    with its simultaneous bit flip.
\end{theorem}
\begin{proof}\leanok
    \uses{thm:peps_czx_pull_through,thm:peps_region_physical_maps}
    Apply the one-site pull-through identity inside the finite contraction.
    Reindex vertex--edge incidences by edge--endpoint incidences. Each
    internal bond contributes the square of its controlled phase, hence
    one; each crossing bond contributes its phase once. Reindex the
    incident-bond sum by simultaneous bit flips.
\end{proof}

\begin{definition}[Oriented rectangle boundary labels]%
    \label{def:czx_rectangle_boundary_labels}
    \lean{TNLean.PEPS.czxRectangleBoundaryLabels,
        TNLean.PEPS.czxRectangleCoordinatesEquiv,
        TNLean.PEPS.czxRectangleCoordinatesEquiv_apply,
        TNLean.PEPS.czxRectangleEffectiveBoundaryConfig,
        TNLean.PEPS.czxRectangleBoundaryMatrix}
    \leanok
    \uses{def:czx_rectangle_perimeter,def:peps_czx_boundary_embedding,
        def:peps_open_region_contraction}
    Let $E_R\mu$ be the clockwise label sequence, with bottom and left
    pairs reversed. This is a bijective change of coordinates. For
    $c\in\{0,1\}^P$, set $\iota(c)_i=(c_i,c_{i+1})$, define
    $\mu(c)=E_R^{-1}\iota(c)$, and let $F_R|c\rangle=A_R|\mu(c)\rangle$.
\end{definition}

\begin{theorem}[The contraction imposes the cyclic plaquette constraint]%
    \label{thm:czx_rectangle_boundary_support}
    \lean{TNLean.PEPS.czxRectangleBoundaryLabels_mem_range_of_openRegionWeight_ne_zero}
    \leanok
    \uses{def:czx_rectangle_boundary_labels}
    If a coefficient of $A_R|\mu\rangle$ is nonzero, then
    $E_R\mu=\iota(c)$ for some $c\in\{0,1\}^P$.
\end{theorem}
\begin{proof}\leanok
    \uses{thm:czx_rectangle_perimeter,thm:czx_open_region_coefficients}
    Extract the nonzero incident-bond assignment. On each straight side,
    the shared internal bond equates the adjacent physical qubits. At each
    corner, the equality is within one site tensor. These eight cases give
    the cyclic adjacent-pair constraint, including the final wraparound.
\end{proof}

\begin{definition}[Plaquette-bit witnesses]%
    \label{def:czx_plaquette_witness}
    \lean{TNLean.PEPS.czxPlaquettePhysical,TNLean.PEPS.czxPlaquetteBondConfig}
    \leanok
    \uses{def:peps_czx_tensor}
    Given bits $q$ at torus plaquette corners, assign site $(x,y)$ the qubits
    $(q_{x,y+1},q_{x+1,y+1},q_{x+1,y},q_{x,y})$ and give each native bond
    the corresponding adjacent pair.
\end{definition}

\begin{theorem}[Exact witnesses and a faithful effective boundary]%
    \label{thm:czx_rectangle_boundary_faithful}
    \lean{TNLean.PEPS.czxPlaquetteBondConfig_right,
        TNLean.PEPS.czxPlaquetteBondConfig_up,
        TNLean.PEPS.component_czxPlaquettePhysical,
        TNLean.PEPS.openRegionWeight_czxPlaquettePhysical,
        TNLean.PEPS.czxRectangleBoundaryLabels_plaquette,
        TNLean.PEPS.czxRectangleEffectiveBoundaryConfig_injective,
        TNLean.PEPS.exists_czxRectangleBoundaryMatrix_selector,
        TNLean.PEPS.czxRectangleBoundaryMatrix_mulVec_injective}
    \leanok
    \uses{def:czx_rectangle_boundary_labels,def:czx_plaquette_witness}
    Each local tensor and each actual region witness coefficient equals
    one. For every $c$ there is a physical configuration $\sigma_c$ with
    $F_R(\sigma_c,c')=\delta_{c,c'}$. In particular, $F_R$ is injective.
\end{theorem}
\begin{proof}\leanok
    \uses{thm:czx_rectangle_perimeter,thm:czx_open_region_coefficients}
    Extend arbitrary boundary bits from the distinct perimeter corners to
    torus plaquette bits, assigning zero elsewhere. The resulting physical
    configuration has coefficient one in the selected column and zero in
    all other columns by disjointness.
\end{proof}

\begin{theorem}[Exact physical image and its dimension]%
    \label{thm:czx_rectangle_physical_image}
    \lean{TNLean.PEPS.range_openRegionMap_czxRectangle,
        TNLean.PEPS.finrank_range_openRegionMap_czxRectangle}
    \leanok
    \uses{def:czx_rectangle_boundary_labels}
    The image of the actual region map is exactly the span of the columns
    of $F_R$, and its dimension is $2^P$.
\end{theorem}
\begin{proof}\leanok
    \uses{thm:czx_rectangle_boundary_support,thm:czx_rectangle_boundary_faithful}
    Every nonzero native boundary column is an effective column; every
    effective column is a native one. The selectors prove independence,
    and there are $2^P$ effective spin configurations.
\end{proof}

\begin{theorem}[Physical action on the effective plaquette spins]%
    \label{thm:czx_rectangle_effective_symmetry}
    \lean{TNLean.PEPS.regionPhysicalMap_czxRectangleBoundaryMatrix_column}
    \leanok
    \uses{def:czx_rectangle_boundary_labels}
    On the faithful physical boundary coordinates,
    \[
        U_RF_R|c\rangle
        =(-1)^{\sum_i c_ic_{i+1}}F_R|\bar c\rangle.
    \]
\end{theorem}
\begin{proof}\leanok
    \uses{thm:czx_region_physical_pullthrough,thm:czx_rectangle_boundary_faithful}
    Reorder crossing bonds clockwise. Pair reversal preserves the controlled
    phase and commutes with bit flip, giving one controlled phase for each
    adjacent pair of effective spins.
\end{proof}

\begin{theorem}[The printed matrix product operator on the physical rectangle image]%
    \label{thm:czx_rectangle_review_mpo}
    \lean{TNLean.PEPS.regionPhysicalProductMatrix_czxRectangleBoundaryMatrix_review}
    \uses{def:czx_rectangle_boundary_labels,def:asymex_czx_review_tensor}
    The review's printed CZX matrix product operator $O(A)$ satisfies
    $U_RF_R=F_RO(A)$ on the actual rectangle image.
\end{theorem}
\begin{proof}
    \uses{thm:czx_rectangle_effective_symmetry,thm:czx_rectangle_perimeter,
        thm:asymex_czx_kernel,thm:asymex_czx_review_kernel}
    The printed matrices give $D_PX^{\otimes P}=(-1)^PX^{\otimes P}D_P$.
    The rectangle perimeter $P$ is even, so this sign is one. The normalized
    bra-ket display is not substituted for those printed matrices.
\end{proof}

The identification above is a linear identification of the actual
open-region image. No formula for the normalized reduced density matrix,
its flatness, or its entropy is asserted here. Nonrectangular regions
remain outside the perimeter construction.
